\chapter{Coupled Fields}

\section{General}
According to the {\sf pdeList} we will introduce a {\sf couplingList}. Within this
list we will differ between iterative coupling ({\sf iterative}) and direct
coupling ({\sf direct}). For iterative coupling the method can be {\sf RHS} 
(for right hand side coupling) and {\sf Matrix} for a matrix coupling algorithm.
Then, for each involved pair of coupling PDEs, we have
to specify for each PDE the coupling parameters. 
The coupling is specified by the coupling input {\sf quantity}, the
interface {\sf type} - which can be {\sf node}(Coupling),
{\sf surface}(Coupling) or {\sf region}(Coupling) -  and the {\sf name} of
the interface, which has to correspond with the according interface
{\sf type}. At the end, we have to define
the stopping criteria and the maximum number of iterations
within the section {\sf nonlLnear}. Therefore, the general structure 
looks as follows


\begin{verbatim}
<couplingList>
  <iterative>
    <Pde1Pde2 method="RHS">
      <Pde1>
        <coupling quantity="..." type="..." name="..."/>
      </Pde1>
      <Pde2>
        <coupling quantity="..." type="..." name="..."/>
        <coupling quantity="..." type="..." name="..."/>
      </Pde2>
    </Pde1Pd2>

    <Pde3Pde4 method="Matrix">
      <Pde3>
        <coupling quantity="..." type="..." name="..."/>
        <coupling quantity="..." type="..." name="..."/>
      </Pde3>
      <Pde4>
        <coupling quantity="..." type="..." name="..."/>
      </Pde4>
    </Pde4Pde4>

    <nonLinear logging="no">
      <stopCrit value="1e-3" quantity="..." l2Norm="..."/>
      <stopCrit value="1e-2" quantity="..." l2Norm="..."/>
      <stopCrit value="1e-5" quantity="..." l2Norm="..."/>
      <maxNumIters> 20 </maxNumIters>
    </nonLinear>
  </iterative>
</couplingList>
\end{verbatim}

\noindent
For the stopping criteria, we have currently the following names:

\begin{verbatim}
    <stopCrit value="1e-2" quantity="mechDisplacement" l2Norm="rel"/>
    <stopCrit value="1e-2" quantity="smoothDisplacement" l2Norm="rel"/>
    <stopCrit value="1e-2" quantity="elecForceVWP" l2Norm="rel"/>
    <stopCrit value="1e-2" quantity="magForceLorentz" l2Norm="rel"/>
    <stopCrit value="1e-2" quantity="acousticForce" l2Norm="rel"/>

\end{verbatim}
Instead of using a relative accuracy $\epsilon^{rel}$, i.e. $\|u_{k+1}-u_k\|_2 < \epsilon^{rel} \|u_{k+1}\|_2$ you can as well
truncate the iterations by choosing an absolute accuracy $\epsilon^{abs}$, i.e. $\|u_{k+1}-u_k\|_2 < \epsilon^{abs}$ setting {\sf
l2Norm="abs"}. Note that the specification of the stopping criterias
is optional: The coupling is considered as converged if either {\sf
  maxNumIters} iterations have been performed or if {\bf all} stopping
criterias were met.


\section{Specifications within the single fields}
Currently we have definitions for electrostatics, mechanics, magnetics and acoustics.

\begin{table}[htb]
\begin{center}
\caption{Electrostatics}
\bigskip
\begin{tabular}{l|l}
input coupling & type \\\hline \hline
nodeCoupling & mechDisplacement \\
regionCoupling & smoothDisplacement 
\end{tabular}
\end{center}
\end{table}

\begin{table}[htb]
\begin{center}
\caption{Mechanics}
\bigskip
\begin{tabular}{l|l}
input coupling    & type \\\hline \hline
nodeCoupling      & elecForceVWP \\
nodeCoupling      & magForceVWP \\
regionCoupling    & magForceLorentz \\
interfaceCoupling & acousticForce 
\end{tabular}
\end{center}
\end{table}

\begin{table}[htb]
\begin{center}
\caption{Magnetics}
\bigskip
\begin{tabular}{l|l}
input coupling    & type \\\hline \hline
nodeCoupling & mechDisplacement \\
regionCoupling & smoothDisplacement 
\end{tabular}
\end{center}
\end{table}

\begin{table}[htb]
\begin{center}
\caption{Acoustics}
\bigskip
\begin{tabular}{l|l}
input coupling    & type \\\hline \hline
nodeCoupling      & mechForce \\
\end{tabular}
\end{center}
\end{table}

\section{ElecMech}
The two involved fields are the electrostatic and the mechanical field.
Often we need a additional PDE for the smoothing of the mesh (moving/deforming 
body problem).

\bigskip
\noindent
{\it Example:}

\begin{verbatim}
<couplingList>
  <iterative>
    <elecMech method="RHS">
      <electrostatic>
        <coupling quantity="smoothDisplacement" type="region" name="air1"/>
        <coupling quantity="smoothDisplacement" type="region" name="air2"/>
      </electrostatic>
      <mechanic>
        <coupling quantity="elecForceVWP" type="node" name="electrode1"/>
      </mechanic>
      <smooth>
        <coupling quantity="mechDisplacement" type="node" name="electrode1"/>
      </smooth>
    </elecMech>
    <nonLinear logging="no">
      <stopCrit value="1e-2" quantity="mechDisplacement" l2Norm="rel"/>
      <stopCrit value="1e-3" quantity="elecForceVWP" l2Norm="abs"/>
      <maxNumIters> 20 </maxNumIters>
    </nonLinear>
  </iterative>>
</couplingList>
\end{verbatim}

\section{MagMech}
The two involved fields are the magnetic and the mechanical field.
Often we need a additional PDE for the smoothing of the mesh (moving/deforming 
body problem).

\bigskip
\noindent
{\it Example:}

\begin{verbatim}
<couplingList>
  <iterative>
    <magMech method="RHS">
      <magnetic>
        <coupling quantity="smoothDisplacement" type="region" name="air1"/>
        <coupling quantity="smoothDisplacement" type="region" name="air2"/>
      </magnetic>
      <mechanic>
        <coupling quantity="magForceLorentz" type="region" name="plate"/>
      </mechanic>
      <smooth>
        <coupling quantity="mechDisplacement" type="node" name="surf1"/>
      </smooth>
      <nonLinear logging="no">
        <stopCrit value="1e-2" quantity="mechDisplacement" l2Norm="rel"/>
        <stopCrit value="1e-3" quantity="magForceLorentz" l2Norm="abs"/>
      <maxNumIters> 20 </maxNumIters>
    </nonLinear>
  </iterative>
</couplingList>
\end{verbatim}


\section{AcouMech}
Here we differ between an indirect and a direct coupling. Currently we have
just implemented the iterative coupling. Note that the coupling
between these two field can only be done via {\sf interface} coupling,
as the element normals are needed for calculating the coupling forces.

\bigskip
\noindent
{\it Example:}

\begin{verbatim}
<couplingList>
  <iterative>
    <acouMech method="RHS">
      <acoustic>
        <coupling quantity="mechForce" type="interface" name="surf1"/>
      </acoustic>
      <mechanic>
        <coupling quantity="acousticForce" type="interface" name="surf1"/>
      </mechanic>
    </acouMech>
    <nonLinear logging="no">
      <stopCrit value="1e-2" quantity="mechDisplacement" l2Norm="rel"/>
      <maxNumIters> 20 </maxNumIters>
    </nonLinear>
  </iterativeAcouMech>
</couplingList>
\end{verbatim}
